\begin{document}
\maketitle
\vspace{-0.7cm}

\begin{keybox}[What this document is, and is not]
This note collects, in one derivation, why the two TFPT inputs leave \emph{no free dimensionless
compiler dial beyond the anchor structure and $\pi$} (the one dimensionful scale $v_{\mathrm{geo}}$ and
the continuous transfer physics $F_{\mathrm{transfer}}$ remain explicitly typed, not derived), and how
the entire integer skeleton, the seam constant $\cthree=\tfrac1{8\pi}$ and a built-in cyclic
structure all flow from one boundary pair. \textbf{Two layers, kept strictly apart:}
\begin{itemize}[leftmargin=1.3em]
\item \textbf{Structural core \mE} (\S\ref{sec:core}--\S\ref{sec:alpha}): exact, machine-checked
identities ($\mathtt{v53}$--$\mathtt{v56}$, plus $\mathtt{v6/\vref{v8}/v3}$). These stand on their own,
independently of any cosmology.
\item \textbf{One interpretation \mC} (\S\ref{sec:cycle}): the \emph{cyclic self-reproduction} reading
--- collapse $\to$ horizon$=$seam $\to$ next cycle. This is a physical hypothesis, \emph{not} derived and
\emph{not} machine-checkable; it is offered because it gives the structural fixed point a physical
selection mechanism. It is flagged \mC\ throughout.
\end{itemize}
Nothing in the \mC\ layer is sold as proven; nothing in the \mE\ layer depends on the \mC\ layer.
\end{keybox}
\markerkey

\tableofcontents

% ---------------------------------------------------------------------
%  tfpt_docset.tex -- THE canonical TFPT document-set box.
%  Single source of truth for the document list; \input by the
%  introduction and every paper. Uses only universally available
%  macros (keybox + plain LaTeX math) so it compiles in every doc.
%  Each document adds its own "You are reading ..." line directly
%  after the \input.  Numbering is aligned with the file names:
%  document N == tfpt_N_*, the introduction is the entry point, and
%  R/H/O are the appendix-style companions.
% ---------------------------------------------------------------------
\begin{navbox}[The TFPT document set --- what is treated where]
\textbf{Plain language:} TFPT (Topological Fixed-Point Theory) is a small discrete compiler. Two
inputs --- the seam constant $c_3=\tfrac1{8\pi}$ and the carrier rank $g_{\mathrm{car}}=5$ --- build
$D_5\oplus A_3+\mu_4\Rightarrow E_8$ and read off the Standard-Model skeleton, the constants and the
scale grammar. The series is \textbf{the introduction plus five numbered papers and four companions}
(plus one standalone working note), best read in order:
\begin{enumerate}[leftmargin=2.0em,itemsep=1pt,label=\textbf{\arabic*.}]
\setcounter{enumi}{-1}
\item \texttt{introduction} --- reading guide, compiler closure, predictions, the dependency DAG and the
      single proof ledger (\texttt{verification/status\_ledger.csv}). \emph{(entry point)}
\item \texttt{tfpt\_1\_architecture\_e8} --- the two axioms, the derivation map, the EM fixed point
      $\alpha^{-1}$, the $D_5\oplus A_3+\mu_4\Rightarrow E_8$ construction, the QBL theorem chain.
\item \texttt{tfpt\_2\_standard\_model} --- the SM in one $\varphi_0$-ladder formula, flavor from
      parabolic transport (sheet diamond, dual anchor, $Q_+$ cohomology), the worked closures.
\item \texttt{tfpt\_3\_e8\_audit\_bootstrap} --- the seven $E_8$ slices as an audit raster, the cascade
      bridge, the M\"obius bootstrap.
\item \texttt{tfpt\_4\_frontier} --- honest status of $\eta_B$, $m_p/m_e$, Koide, dark matter, full QG.
\item \texttt{tfpt\_5\_redteam} --- the adversarial audit: declared attacks, what survives, what each
      reduces to, and the kill tests.
\end{enumerate}
\noindent\textbf{Companions:} \textbf{R.} \texttt{tfpt\_research\_contracts} --- the open interfaces
$(v_{\mathrm{geo}}, G_{\mathrm{net}}, F_{\mathrm{transfer}})$ as numbered contracts;
\textbf{H.} \texttt{tfpt\_horizon\_readouts} --- Appendix~H (unit-system reframe): one seam constant as
the universal horizon thermal code, the Nariai/anchor surface, the resummed seam clock;
\textbf{O.} \texttt{origin\_theory} --- why no free fundamental number remains (exact core $+$ one
honestly-typed cyclic interpretation);
\textbf{S.} \texttt{tfpt\_safeguards} --- the verification discipline: every mechanism that defends a claim
against coincidence and numerology (status calculus, no-free-pattern, the over-determination map, the
firewall, frozen predictions $+$ null model, the independent Wolfram and Lean paths, the red team).\\[2pt]
\noindent\textbf{Working notes:} \textbf{N1.} \texttt{note\_e8\_gaussian\_code} --- the Gaussian code
bridge: $E_8$ over $\mathbb Z[i]$ via Construction~A over the extended Hamming code $[8,4,4]$, the
canonical four-bit quotient $L/(1{+}i)L\cong\mathbb F_2^4$, and the $G_{31}$ quartic companion (exact
algebra, machine-certified).
\end{navbox}

\noindent\textbf{You are reading \textbf{companion~O} (Origin Theory).}

\input{tex-artefacts/tfpt_status}
\TFPTbadges
\TFPTclaimstack[anchor]

\bigskip

% =====================================================================
\section{The whole integer skeleton from one pair $(\gcar,\Nfam)=(5,3)$}\label{sec:core}
% =====================================================================
TFPT runs on two operational inputs: the seam constant $\cthree=\tfrac1{8\pi}$ (P1) and the carrier rank
$\gcar=5$ (P2). The boundary $\mu_4=\{1,i,-1,-i\}$ gives the four-punctured line $\mathbb P^1\setminus\mu_4$
with $H_1=\Z^3$, hence the family count $\Nfam=3$ ($A_3$). We show the rest of the discrete data is
generated by the single pair $(\gcar,\Nfam)=(5,3)$.
\par\smallskip
\noindent\emph{One object, two canonical invariants.} The pair $(\gcar,\Nfam)$ is read off the
\emph{same} four-point divisor $D=\mu_4$ (degree $4$) by the two canonical functors on $\PP^1$: the cycle
side $\operatorname{rank}H_1(\PP^1\setminus\mu_4)=\deg D-1=3=\Nfam$ and the function side (Riemann--Roch)
$h^0(\PP^1,\mathcal O(\mu_4))=\deg D+1=5=\gcar$, the latter matched by $\operatorname{rank}(D_5)=5$ (the
carrier is the $\mathfrak{so}(10)$ half-spinor). Reading the carrier as this meromorphic seam-mode space
(at most simple poles on $\mu_4$) is geometrically natural \mE\ in its arithmetic, though the physical
identification stays \mC\ and does not displace the over-determined $\gcar=5$ (Pascal/bootstrap/Lean)
\veri{v189\_riemann\_roch\_carrier.py}.
\par\smallskip
\noindent\emph{And the function side generates $D_5$ itself, not just $\gcar$.} The even Clifford algebra of
the $5$-dimensional mode space closes the chain $\mu_4\to\gcar\to D_5$: $\Lambda^{\mathrm{even}}(C_{\mathrm{car}})
=\binom50+\binom52+\binom54=1+10+5=16=2^{\gcar-1}=\dim S^+$, which is exactly the $\mathfrak{so}(10)=D_5$
half-spinor. So the carrier is not merely ``rank $5$ like $D_5$'' --- it \emph{is} the Clifford spinor of the
Riemann--Roch mode space \mE, with the physical identification again \mC\ \veri{v197\_rr\_carrier\_clifford\_d5.py}.

\begin{keybox}[The $(5,3)$ generator and the Pythagorean mass volume \mE\ (\veri{v53\_compiler\_core.py})]
By sum, difference and halving,
\[
  \operatorname{rank}E_8=\gcar{+}\Nfam=8,\qquad |\Z_2|=\gcar{-}\Nfam=2,\qquad |\mu_4|=\tfrac{\gcar+\Nfam}{2}=4,
\]
and $(\Nfam,|\mu_4|,\gcar)=(3,4,5)$ is \emph{the} Pythagorean triple. This is load-bearing: the
grand-mass-volume exponent $\Delta_Y=\gcar^2=25$ ($\det M_{\mathrm{SM}}\sim(\phiz)^{25}$) splits as a
\emph{difference of squares},
\[
  \boxed{\ \Delta_Y=\gcar^2=\underbrace{\Nfam^2}_{\text{down}=9}
  +\underbrace{(\gcar{-}\Nfam)(\gcar{+}\Nfam)}_{=\,|\Z_2|\cdot\operatorname{rank}E_8=\dim S^+=16}
  =9+16=25\ }
\]
(the $K$-row sums $(6,9,10)$: down carries $\Nfam^2$, up$+$lepton carry $|\Z_2|\!\cdot\!\operatorname{rank}E_8$).
The \emph{anchor} $a=(1,1,2)$ is the unique $3$-multiset whose elementary symmetric functions are the
compiler atoms $(|\mu_4|,\gcar,|\Z_2|)=(4,5,2)$; equivalently $\chi_a(t)=(t{-}1)^2(t{-}2)=t^3-|\mu_4|t^2+\gcar t-|\Z_2|$.
Sharper (\veri{v491\_p2\_partition\_corollary.py}): only $e_1{=}4$ is needed --- three positive integers
summing to $4$ already force $\{1,1,2\}$, so $\gcar{=}5$ and $|\Z_2|{=}2$ are corollaries of the four
marks given the weight typing; \texttt{AX.P2.01} stays the declared axiom (the residual is the typing).
And the typing is itself hardened (\veri{v499\_p2\_weights\_deligne\_bg.py}): Deligne $+$
Birkhoff--Grothendieck $+$ Mehta--Seshadri stability derive integrality, positivity ($h^0{=}0$) and the
sum $4$ from the four marks, rank $3$, the cusp class and (U) --- the \mC\ residual is the module
identification plus (U) itself.
\end{keybox}

The two glue discriminant norms are likewise $(\gcar,\Nfam)$ over the glue index: $q(D_5)=\gcar/|\mu_4|=\tfrac54$,
$q(A_3)=\Nfam/|\mu_4|=\tfrac34$, so their sum is the $E_8$ root norm $2$, their difference is the lepton
transport value $\delta=|\Z_2|/|\mu_4|=\tfrac12$, their product is $\dim\mathfrak{su}(4)/\dim S^+=\tfrac{15}{16}$
and their ratio is $\gcar/\Nfam=\tfrac53$ \veri{v51\_boundary\_half\_step.py}. The single transcendental is
$\pi$; both $\pi$-coefficients are $2\pi$ times a compiler integer.

\begin{keybox}[The icosahedral bedrock: \emph{why} the atoms are $\{2,3,5\}$ \mE\ (\veri{v219\_icosahedral\_mckay.py})]
The pair $(\gcar,\Nfam)=(5,3)$ and the sheet $|\Z_2|=2$ are not three free small integers: $\{2,3,5\}$ is
forced as soon as one fixes $E_8$. By the \textbf{McKay correspondence} the finite subgroups of $SU(2)$ are
classified by the affine ADE diagrams, and $E_8$ is the exceptional \emph{maximum} of that spherical tower
($2T{\to}\hat E_6$, $2O{\to}\hat E_7$, $2I{\to}\hat E_8$). Its branch data is the \emph{icosahedron}: the
binary icosahedral group $2I$ (order $120{=}|R^+(E_8)|$) has element orders $\{1,2,3,4,5,6,10\}$ --- it
literally \emph{contains} the $2,3,5$ rotation-axis orders --- and its nine irreducible representations have
dimensions $\{1,2,2,3,3,4,4,5,6\}$, the Kac marks of $\hat E_8$, with
\[
  \sum_i d_i = 30 = h(E_8) = |\Z_2|\cdot\Nfam\cdot\gcar = 2\cdot3\cdot5,\qquad
  \sum_i d_i^2 = 120 = |R^+(E_8)| .
\]
The McKay graph is \emph{built from the group} (the natural $2$-dim rep gives the affine $E_8$ adjacency,
top eigenvalue $2$, with marks $=$ degrees), so the recurring $\{2,3,5\}$ signature has a single bedrock:
$E_8$ \emph{is} the icosahedron. This is a \emph{backward certificate} of the already-closed $E_8$ \mE; the
reading ``it explains the $(2,3,5)$ choice'' is \mC; the raw-seam origin (without first positing $E_8$) stays
the \mO\ P2 interface. The same exceptional geometry has two metric readings via complex multiplication
(\veri{v222\_cm\_norm\_duality.py}): on the \emph{square} seam ($j{=}1728$, Gaussian $\Z[i]$) the EM index is
the Gaussian norm $N(\gcar{+}i|\mu_4|)=5^2{+}4^2=41=10b_1$, while on the \emph{hexagonal} partner ($j{=}0$,
Eisenstein $\Z[\omega]$) the scalaron exponent is $N(3{+}2\omega)=7$ --- the carrier split $(3,2)$ generates
$(5,6,7,13)$ by sum/product/Eisenstein/Gauss norm, the rings rigid. The same two rings give the atoms a
\emph{splitting trichotomy} (\veri{v416\_atom\_trichotomy.py}): in $\Z[i]$ (the square seam) $2{=}|\Z_2|$
ramifies, $3{=}\Nfam$ is inert and $5{=}\gcar$ \emph{splits} as $(2{+}i)(2{-}i)$, while in $\Z[\omega]$ (the
hexagonal partner) $3$ ramifies and $2,5$ are inert --- so the ramified prime of each ring \emph{is} its own
atom, and the three quadratic facets $\mathbb Q(i),\mathbb Q(\sqrt{-3}),\mathbb Q(\sqrt5)$ of the order-$30$
field $K=\mathbb Q(i,\sqrt{-3},\sqrt5)$ (\vref{v403}) ramify over exactly one atom each ($2,3,5$; product
$30=h(E_8)$), two imaginary (CM/dynamic) plus one real (RM/static). The arithmetic is \mE, the
seam/flavor/golden reading \mC. The same three rings carry the carrier split $(3,2)$ as one
determinant triple (\veri{v418\_cyclotomic\_norm\_triple.py}): $\det(3I+2\,\mathrm{Comp}(\Phi_n))=(7,13,55)$
for $n=3,4,5$ over $\Z[\omega],\Z[i],\Z[\zeta_5]$ --- the scalaron, $\Delta_Q$ and the quark numerator
$55=\gcar\|\operatorname{Pl}(K)\|_1$; the carrier-$5$ clock missing from the $3$-dim flavor space is the
$4{\times}4$ $\Phi_5$-companion (with the golden $\varphi$ as its real-part data). The anchor $(5,4)$ fails
the triple in the carrier ring ($461$, prime), and honestly $(3,2)$ is the \emph{forced} split (\vref{v14}),
not the unique clean one (the scan also finds $(2,1){\to}(3,5,11)$, a secondary rung). The sharper
discriminator is \emph{cross-sector}: only $(3,2)$ lights up three different physics sectors --- gravity
($7$), flavor ($13$) and masses ($55$) --- whereas $(2,1){\to}(3,5,11)$ is all compiler-core
($\Nfam,\gcar,\|\operatorname{Pl}(K)\|_1$); the two rungs read as input $(2,1)\to$ output $(3,2)$. \mE/\mC.
And the seam $\mu_4$
clock is the
order-$4$ character of the $E_8$ Coxeter cycle itself: $(\Z/30)^\times$ are the $E_8$ exponents, with the
order-$4$ generator $7$ acting as a literal $\mu_4$ clock on the eight live phases
(\veri{v223\_coxeter\_totative\_clock.py}).
\par\smallskip\noindent\emph{The seam $\mu_4$ is the carrier's Galois group
(\veri{v419\_seam\_galois\_carrier.py}).} The deeper reason the seam sits on the Galois side, not on the
cyclic dial, is that $h(E_8)=30$ is \emph{squarefree}: $\Z/30$ has no order-$4$ element, so $\mu_4$ cannot
be a rotation hand. By CRT $(\Z/30)^\times=(\Z/2)^\times\times(\Z/3)^\times\times(\Z/5)^\times=\mu_4\times\Z_2$
(order $\varphi(30)=8=\operatorname{rank}E_8$), and the $\mu_4$ factor \emph{is}
$(\Z/5)^\times=\operatorname{Gal}(\mathbb Q(\zeta_5))$ --- the carrier prime $5$. So the seam $\mu_4$ is the
carrier pentagon's automorphism group: the Frobenius $\zeta_5\mapsto\zeta_5^2$ is the explicit order-$4$
operator $G$ with $G\,C_5\,G^{-1}=C_5^2$, $G^4=I$ on the carrier clock $C_5$ (\vref{v418}). This is the
positive content of \texttt{RES.COXETER.SYMMETRY.01} (\vref{v409}): the $2^2=4$ is the order of
$(\Z/5)^\times$, and the seam ``carries no rate'' because it is an \emph{automorphism}, not a hand. So
$\mathbb Q(\zeta_{30})$ carries both $E_8$ invariants --- the cyclic hands ($h=30$) and the Galois gears
($\operatorname{rank}=8=\mu_4\times\Z_2$).
\begin{center}
\includegraphics[width=0.92\linewidth]{coxeter_galois.pdf}
\end{center}
\par\smallskip\noindent\emph{The Galois gearbox IS the $(E_8)_1$ glue
(\veri{v422\_galois\_net\_bridge.py}).} This bridges the Galois side directly to $G_{\mathrm{net}}$: the seam
$\mu_4$ is not merely \emph{a} cyclic $\Z/4$, it is the \emph{same} one as the holomorphic-net glue. The
carrier discriminants $\operatorname{disc}(A_3)=\operatorname{disc}(D_5)=\Z/4$ are cyclic (one Smith
invariant factor $4$), and $D_5$ is cyclic precisely because its rank $5=\gcar$ is \emph{odd} ($D_n$ is
$\Z/4$ for $n$ odd, $\Z_2\times\Z_2$ for $n$ even); the carrier-Galois
$\operatorname{Gal}(\mathbb Q(\zeta_5))=(\Z/5)^\times$ is cyclic $\Z/4$ for the same reason ($5$ prime); and
in $\operatorname{disc}(D_5\oplus A_3)=\Z_4\times\Z_4$ the simple current $(1,1)$ generates the cyclic $\Z/4$
whose Lagrangian quotient $16/4^2=1$ \emph{is} $(E_8)_1$ (\vref{v89}/\vref{v154}). The cyclic-vs-Klein
distinction is load-bearing: the order-$2$ halfway glue gives $D_8=SO(16)$ ($\det4$), not $E_8$ ($\det1$),
so the carrier-Galois (cyclic, $5$ prime) matches $E_8$, never the $D_8$ Klein. So the order-$30$
\emph{gearbox} $(\Z/30)^\times$ is the simple-current glue of the seam net --- one $\Z/4$ from carrier
arithmetic through the lattice to the seam (\mE; the canonical identification \mC). This is the
\emph{four-fold $\mu_4$ identity} --- deck/divisor $=$ Galois $=$ discriminant $=$ simple-current, one
cyclic $\Z/4$ read four ways --- and, by the over-determination accounting (\vref{v428}), it counts as
\emph{compression} (a unification that makes the carrier less axiomatic), not four independent witnesses
that multiply.
Finally, P2 itself reads as a Riemann--Roch index: a degree-$4$
seam divisor on $\PP^1$ forces $h^0{=}5{=}\gcar$ and $\operatorname{rank}H_1{=}3{=}\Nfam$, with even Clifford
closure $\Lambda^{\mathrm{even}}(\mathbb C^5)=16=\dim S^+$ (\veri{v228\_rr\_index\_gate.py}, \mE/\mC/\mO,
conditional on the seam producing the degree-$4$ divisor $=$ \texttt{QGEO.SYM.01}). And the same
icosahedron gives the seam a \emph{geometric} home (\veri{v232\_e8\_kleinian\_seam.py}): on the du~Val side
of McKay, the minimal resolution of the Kleinian singularity $\mathbb C^2/2I$ has an exceptional locus of
\emph{eight} rational curves glued in the $E_8$ pattern (their negated intersection form is the $E_8$ Cartan,
$\det1$; each a $(-2)$-curve), and its link is the Poincar\'e homology $3$-sphere $S^3/2I$ ($2I$ is perfect,
so $H_1=0$, $\pi_1=2I$ of order $120=|R^+(E_8)|$). The eight curves are $\operatorname{rank}E_8=\gcar+\Nfam$
--- a \emph{fourth} reading of the ``$8$'' as the eight seam components. So the raw seam ($\PP^1$ with marks)
is canonically one of the $E_8$ du~Val exceptional curves: \texttt{QGEO.REALIZE.01} acquires an
algebraic-geometric model \mC, though it still presupposes $E_8$ (the bedrock stays \mO).
\par\smallskip\noindent\emph{The capstone: one singularity generates the skeleton
(\veri{v236\_brieskorn\_capstone.py}).} All of the above is the Milnor fibration of a \emph{single}
singularity --- the $(2,3,5)$ Brieskorn singularity $x^2+y^3+z^5$, whose three exponents \emph{are} the atoms
$(|\Z_2|,\Nfam,\gcar)$. Its Milnor number $\mu=(2{-}1)(3{-}1)(5{-}1)=1{\cdot}2{\cdot}4=8=\operatorname{rank}E_8$
is a \emph{fifth} independent origin of the ``$8$'' (after geometry, lattice, gravity and $\varphi(30)$);
its Milnor monodromy has eigenvalues the eight primitive $30$th roots $e^{2\pi i m/30}$ with $m$ the $E_8$
exponents (characteristic polynomial $\Phi_{30}$), i.e.\ it \emph{is} the order-$30$ $E_8$ Coxeter element
(\S\ref{sec:coxeter}); its Milnor lattice is $E_8$ and its link is the Poincar\'e sphere. Both clocks are
facets of this one monodromy: the $\mu_3$ family/triality (the CP phase, \veri{v233\_cp\_triality\_phase.py})
is the order-$3$ power $h^{10}$ in $\langle h\rangle=\Z/30$, while the $\mu_4$ seam clock
(\veri{v223\_coxeter\_totative\_clock.py}) is the order-$4$ Galois automorphism of $\mathbb Q(\zeta_{30})$. So the
recurring $\{2,3,5\}$ is not many numbers but \emph{one} icosahedral singularity read several ways \mE.
\end{keybox}

\par\smallskip\noindent\textbf{The atoms are the network's spectral angles --- the golden ratio is the
$\gcar{=}5$ signature (\veri{v313\_golden\_atoms\_spectral.py}).} The same $(2,3,5)$ has a spectral face. The
affine-$E_8$ adjacency characteristic polynomial factors as $x(x^2{-}4)(x^2{-}1)(x^2{-}x{-}1)(x^2{+}x{-}1)$, so its
non-trivial eigenvalues are exactly $2\cos(\pi/k)$ for $k\in\{|\Z_2|,\Nfam,\gcar\}=\{2,3,5\}$: $2\cos(\pi/2){=}0$,
$2\cos(\pi/3){=}1$, $2\cos(\pi/5){=}\phi$ (the golden ratio, with the second harmonic $2\cos(2\pi/5){=}1/\phi$).
\mE\ So the golden ratio is \emph{specifically} the $\gcar{=}5$ ($5$-fold) signature, and $\gcar{=}5$ gains a new
spectral witness ($\pi/5$) distinct from Pascal and Riemann--Roch. The selection is the icosahedral one:
$1/|\Z_2|{+}1/\Nfam{+}1/\gcar=31/30>1$ is the spherical-triangle condition forcing the \emph{finite}
$2I\Rightarrow E_8$, and its numerator $31=2^{\gcar}{-}1=248/8=1{+}h(E_8)$ is exactly the gap-decoupling margin
numerator (\vref{v63}/\vref{v76}) over the Coxeter $30=h(E_8)$ --- a unification of already-anchored quantities,
not new evidence. Honestly, $\phi$ is a \emph{structural} lattice/network quantity, absent from the dimensionless
SM readouts --- with the sole \mC\ exception that $\cos(\theta_i{=}3\pi/5)=-1/(2\phi)$ on the axion spine branch,
where $\theta_i$ is the regular $\gcar$-gon (pentagon) interior angle (\veri{v429\_axion\_pentagon\_phi.py}); and
as the $E_8$-quasicrystal (Elser--Sloane cut-and-project) ratio it is a concrete candidate substrate
for the rewrite question \mC.

\par\smallskip\noindent\textbf{A minimal substrate that works}
(\veri{v324\_hypergraph\_fiber.py}). The rewrite question can be made concrete with a \emph{fibred} $(2,3,5)$
network. The pure $(2,3,5)$ network gives only
the Coxeter skeleton and the golden $5$-fold angle; the recovery rate $(2/3)^6$ is provably \emph{not} in its
adjacency spectrum (\vref{v312}). But the \emph{smallest} substrate that carries everything is a simple
\emph{product}: the carrier network $T_{\mathrm{net}}=(A{+}2I)/4$ (whose attractor is the Kac marks $=$ the $E_8$
skeleton) fibred by a $3$-node family cusp $T_{\mathrm{cusp}}=\operatorname{diag}((1{-}w)^{2\Nfam})$,
$w\in\{0,\tfrac13,\tfrac23\}$ (whose spectrum $\{1,(2/3)^6,(1/3)^6\}$ supplies the recovery gap). \mE\ The fibred
operator $T_{\mathrm{net}}\otimes T_{\mathrm{cusp}}$ carries \emph{both} at once: its top eigenvalue $1$ has
eigenvector $\mathrm{marks}\otimes(w{=}0)$ (the skeleton $\times$ the democratic family law), and $(2/3)^6$ appears
as a genuine eigenvalue with eigenvector $\mathrm{marks}\otimes(w{=}\tfrac13)$ (network attractor $\times$ cusp
subleading). \mE\ This is the cyclotomic split made dynamical: substrate $=$ carrier $\times$ family $=$ the
$30=\gcar(2\Nfam)=5\times6$ of \vref{v315}, the recovery gap living entirely in the family factor, and the minimal
fibre the $3$-node cusp ($=\Nfam$). Honestly (\vref{v312}) the cusp weight is still \emph{injected}, not derived from
the graph --- so this is a consistent realisation of the predicted product structure, not a derivation of $2/3$ from
a pure rewrite \mC.

\par\smallskip\noindent\textbf{The cusp weight from a minimal rewrite}
(\veri{v327\_hypergraph\_rewrite.py}). One honest step further: the injected $2/3$ is the fixed point of a
\emph{minimal local rule}. A family channel with $\Nfam{=}3$ slots, one absorbing attractor ($w{=}0$) and
$|\Z_2|{=}2$ surviving (the $\Z_2$ pair), has per-step survival operator
$M=\bigl(\begin{smallmatrix}1&0&0\\0&\frac13&\frac13\\0&\frac13&\frac13\end{smallmatrix}\bigr)$ with spectrum
$\{1,\tfrac23,0\}$ --- so the recovery survival $\tfrac23=(\Nfam{-}1)/\Nfam=|\Z_2|/\Nfam$ \emph{emerges from the
rule arity}, and over the order-$6{=}2\Nfam$ hand the rate is $(2/3)^6$. \mE\ Moreover $2/3$ is provably \emph{not}
a root of the affine-$E_8$ characteristic polynomial ($p(2/3)=-3520/19683\ne0$), so the recovery rate cannot be an
adjacency eigenvalue (the \vref{v312} negative, now a proof). \mO\ The residual shrinks: the rewrite derives the
\emph{ratio} $2/3$ from the arity $\{|\Z_2|,\Nfam\}=\{2,3\}$, leaving only the anchor atoms as input --- so the
substrate is the $(2,3,5)$ adjacency (carrier) $\times$ this minimal branching rule (family), both built from
$\{2,3,5\}$ \mC.

\par\smallskip\noindent\textbf{The full spectrum from the unique lazy rule}
(\veri{v486\_transfer\_full\_rule.py}). That rule generates only $\lambda_2$ --- its zero mode persists under
every power, so no iterate reaches the verified \emph{third} transfer mode $(1/3)^6$ (\vref{v54}/\vref{v221}),
which so far had no local-rule origin. \mE\ The gap closes with a \emph{uniqueness} statement: for the symmetric
rule $M(s,h)$ (one absorbing family channel $+$ $\Z_2$ pair) the survival spectrum is $\{s{+}h,\,s{-}h\}$, and
demanding the \emph{physical} pair $\{\tfrac23,\tfrac13\}$ forces uniquely
$(\mathrm{stay},\mathrm{hop},\mathrm{leak})=(\tfrac12,\tfrac16,\tfrac13)=(1/|\Z_2|,\,1/(|\Z_2|\Nfam),\,1/\Nfam)$
--- the \emph{lazy} $\Z_2$-pair walk, every rate an atom expression. Over the same order-$6$ hand it generates
\emph{exactly} $\operatorname{eig}(B^6)=\{(2/3)^6,(1/3)^6\}$, so both decay gaps now have one recursive
generator: $6\ln\tfrac32$ (recovery, \vref{v76}) and $6\ln3$ (subdominant, previously generator-less). \mO\ As
before the arity $\{2,3\}$ is the anchor input; \emph{additionally} the lazy-vs-uniform split is selected by the
physical spectrum, not yet derived (named candidate: stay $=1/|\Z_2|$ as the sheet-symmetric choice of the
$\Z_2$ involution).

\par\smallskip\noindent\textbf{The lazy split forced by the resummed clock}
(\veri{v487\_transfer\_clock\_rungs.py}). The split selection closes --- and onto \emph{established} structure
rather than a new principle. The lazy rule's one-step spectrum is exactly the complete resummed-clock ladder
below the wall (\vref{v124}): $\operatorname{eig}(M)=\{1-n/\Nfam:n=0,1,2\}=\{1,\tfrac23,\tfrac13\}$, while the
uniform rule puts its second pair mode \emph{on} the wall ($0=1-3/3$) --- the degeneracy re-read as ``odd channel
collapsed to the wall''. \mE\ The deck parity \emph{is} the rung index: $[B,\sigma]=0$, the $\sigma$-even pair
mode carries rung $n{=}1$ (survival $\tfrac23$), the $\sigma$-odd mode rung $n{=}2$ (survival $\tfrac13$) ---
each clock step flips the deck parity, the structural home of the parity assignment ($\omega_1$ even-channel,
$\omega_2$ odd-channel) that the empirical comb searches test. \mE\ Forcing: demanding the \emph{faithful} ladder
(no rung missing, no physical mode on the wall --- definitional for a full-spectrum generator), $\Z_2$
equivariance and non-negative rates fixes uniquely both the split $(\tfrac12,\tfrac16)$ \emph{and} the rung
assignment even${}\to{}$1, odd${}\to{}$2 (the swap needs $h=-\tfrac16<0$); the ``sheet-symmetric choice''
candidate is superseded by this proof. \mE\ Corollary: $\omega_1/\omega_2=\mathrm{rate}(2)/\mathrm{rate}(1)
=\log_{3/2}3$ --- the two comb tones are one \vref{v107}/\vref{v124} bend apart. \mO\ What remains open is
exactly what was open before v486: the arity $\{2,3\}$ (anchor input) and the clock's semiclassical derivation
(R1); no independent choice remains at the walk level.

\par\smallskip\noindent\textbf{One coupled local rule}
(\veri{v395\_hypergraph\_coupled.py}). The three hypergraph modules (\veri{v299\_hypergraph\_0d\_autonomous.py},
\veri{v327\_hypergraph\_rewrite.py}, \veri{v324\_hypergraph\_fiber.py}) merge
into a \emph{single} rewrite on a $9{\times}3$ labelled grid: one micro-step $=$ network lazy diffusion on each
cusp column $+$ the \veri{v327\_hypergraph\_rewrite.py} rule $M$ on each node's family vector ($T_{\mathrm{micro}}{=}T_{\mathrm{net}}\otimes
M$, purely local on $27$ cells). \mE\ The joint attractor is marks$\otimes(w{=}0)$; one clock hand ($2\Nfam$ family
micro-steps) carries recovery $(2/3)^6$; \veri{v324\_hypergraph\_fiber.py}'s $T_{\mathrm{net}}\otimes T_{\mathrm{cusp}}$ is the same
substrate in the Galois/cusp readout basis (\veri{v317\_galois\_family.py}); \veri{v299\_hypergraph\_0d\_autonomous.py}'s growth $E_6{\to}E_7{\to}E_8{\to}\hat E_8$ is
unchanged with the fiber attached. \mO\ The mechanism is unified; the irreducible inputs remain the $3$-arm seed
($P_2$) and the analytic seed $\phiz$ (puncture term still \veri{v312\_hypergraph\_substrate.py}).

\par\smallskip\noindent\textbf{The seed at leading order}
(\veri{v396\_phi0\_icosahedral.py}). The tree-level retained seed
$\varphi_0^{\mathrm{tree}}=1/(6\pi)$ is not an independent hypergraph injection: on the icosahedral
$(2,3,5)$ skeleton it equals $F/(4h\pi)=F/(|\Aut|\pi)=(F/(\gcar\Nfam))c_3$ with $F=20$ triangular
hyperedges, $h=30$, $|\Aut|=120$, and $F/h=|\Z_2|/\Nfam=2/3$ (the same ratio as \veri{v327\_hypergraph\_rewrite.py}). \mE\
Gauss--Bonnet consistent ($c_3=1/(|\Z_2|4\pi)$, \veri{v216\_marks\_gauss\_bonnet.py}/\veri{v342\_em\_ward\_heatkernel.py}).
\mE\ The puncture $48c_3^4$ is still analytic, not a rational hypergraph fraction
(\veri{v312\_hypergraph\_substrate.py}); it belongs to the \emph{boundary} engine, not the graph --- named
\texttt{HYP.PHI0.PUNCTURE.01} (\veri{v408\_phi0\_puncture\_heatkernel.py}): the order-$|\mu_4|{=}4$ local
boundary heat-kernel (Seeley--DeWitt) contact term $\Omega_{\mathrm{adm}}c_3^4$ of the $\mu_4$ puncture divisor,
summed over $\Omega_{\mathrm{adm}}{=}\Nfam\dim S^+{=}48$ admissible states ($c_3$ the same Gauss--Bonnet datum,
so $c_3^4$ a fourth-order boundary power; cf.\ \vref{v391}). \mC\ Two engines, two terms: the icosahedral
hypergraph delivers the tree term, the boundary heat-kernel delivers the puncture; its from-first-principles
heat-kernel proof is the tracked \mO\ residual. \mE\ That residual is meanwhile \emph{halved}
(\veri{v483\_pillowcase\_exact\_traces.py}): the geometric/equivariant side is EXACT --- every twisted heat
trace of the flat $\tau{=}i$ pillowcase is a $t$-\emph{independent rational}
($\operatorname{Tr}(\sigma e^{t\Delta})=\operatorname{Tr}(\rho e^{t\Delta})=1$ for all $t$, the Atiyah--Bott
fixed-point counts $4{\times}\tfrac14$ and $2{\times}\tfrac12$; pillowcase contact term exactly $\tfrac12$;
clock-equivariant trace exactly $1$) --- so the $\pi^{-4}$ in $48c_3^4=\tfrac{3}{256\pi^4}$ \emph{cannot}
come from flat orbifold geometry at any order: it must sit in the per-insertion \emph{coupling}
normalisation ($4{=}|\mu_4|$ marks of weight $c_3$ give $c_3^4$, multiplicity $48{=}\Omega_{\mathrm{adm}}$;
the bare $1/(2\pi)$-propagator $4$-cycle differs by the exact rational $\tfrac{3}{16}$). The target narrows
from ``derive the term'' to \emph{one} rule --- ``per-mark insertion weight $=c_3$'' (the same residual
class as the EM-Ward $c_3$-ladder) --- and stays \mO. \mE\ And that rule is itself \emph{not a new
unknown} (\veri{v484\_seam\_contact\_unit.py}): ``$c_3$ per insertion'' is the KMS seam unit
$2\pi=1/(4c_3)$ with $\tfrac14=1/|\mu_4|$ (\vref{v239}) --- one \emph{bare} boundary propagator
($\tfrac1{2\pi}=4c_3$) orbit-averaged over the four marks. On the actual seam circle this is derived for
the finite cycle sector: the bare Green function $G(d)=-\tfrac1\pi\ln|2\sin(d/2)|$ takes \emph{integer}
multiples of $c_3\ln2$ at the $\mu_4$ separations ($G(\pi/2)=-4c_3\ln2$, $G(\pi)=-8c_3\ln2$, mark matrix
$-4c_3\ln2\cdot\mathrm{circ}(0,1,2,1)$, spectrum $\{4,-2,0,-2\}$), so the log-det contact expansion is
\emph{exactly} graded in $c_3$ per insertion ($\operatorname{Tr}C^k=4^k+2(-2)^k$). The grammar is already
suite-wide: the $\Lambda$ prefactor $\tfrac{3}{4\pi^2}=(8\pi)^2 d_{\mathrm{top}}=48c_3^2$ (\vref{v60})
carries the \emph{same} multiplicity $48=\Omega_{\mathrm{adm}}$ at \emph{two} insertions, and the $g{-}2$
seam loop divides by the same $2\pi=1/(4c_3)$. Consequence \mC/\mO: this target and
\texttt{ALPHA.QUILLEN.EXACT.01} \emph{merge} into one remaining analytic step --- the diagonal
(self-energy) zeta-renormalisation at the marks plus the state-multiplicity matching
($48=\Omega_{\mathrm{adm}}$, $41=10b_1$); the finite half is proven, no coefficient is claimed. \mE\ That
remaining step is meanwhile settled at the computable level (\veri{v485\_contact\_diagonal\_closed.py}):
the renormalised diagonal vanishes \emph{exactly} at the KMS seam circumference
($G_{\mathrm{reg}}(0;\ell)=\tfrac1\pi\ln(\ell/2\pi)$, zero iff $\ell=2\pi=1/(4c_3)$), the mark
determinant resums closed-form ($\det(I-uC)=(1-4u)(1+2u)^2$, $u=\varepsilon c_3\ln2$; BFK route,
\vref{v151}), and the $48$/$41$ multiplicities are \emph{one} state set under two response weights (flat
vs $Y^2$/Ginsparg). The merged target's remaining \mO\ is the abstract-seam $\zeta$-det identification
--- a face of \texttt{SEAM.EQUIV.01} alone. \mE\ That face is meanwhile \emph{executed} on the zeta side
(\veri{v974\_alpha\_faces\_computed.py}): the BFK Mayer--Vietoris formula runs on the seam circle cut at
the four $\mu_4$ marks (gluing constant exactly $2^{-N}$, $\det\nolimits'_\zeta(\Delta_{S^1})=\ell^2$,
jump matrix $=$ the inverse mark Green matrix), and at $\ell=2\pi=1/(4c_3)$ it reproduces the \vref{v485}
diagonal/scale channel exactly --- but against the \vref{v484} mark matrix there is \emph{no} global
proportionality: the two kernels sit in \emph{swapped} $\mu_4$ channels (gauge zero mode vs.\ sign rep),
channel-wise exactly matched (same $\mu_4$ Fourier eigenbasis, $|\mathrm{spec}|$ multisets $\{0,2,2,4\}$,
doublet ratio $4/\ln2$ with $c_3$ cancelling). The remaining \mO\ is thereby \emph{sharpened}, not
closed: the channel-swap question plus the abstract-seam identification --- and the swap half is
meanwhile resolved at the \emph{finite} level by $\mu_4$-character grading (deck swap $r\to r{+}2$,
graded $\det'$ difference $c_3$-free $=\log(\ln2/4)$, uniform channel ratio $4/\ln2$;
\veri{v985\_quillen\_channel\_swap.py}), leaving the abstract-seam/continuum identification as the
open residual.

\par\smallskip\noindent\textbf{Static vs.\ dynamic rates: a number-field split}
(\veri{v314\_rate\_translation.py}). The discrete kernel and the dynamic recovery sit in \emph{different number
fields}, which is the precise reason their rates do not freely interconvert. The adjacency spectrum splits as a
rational part $\{0,\pm1,\pm2\}\subset\mathbb Q$ and a golden part $\{\pm\phi,\pm1/\phi\}\subset\mathbb Q(\sqrt5)$,
while the dynamic transfer spectrum $\{1,(2/3)^6,(1/3)^6\}$ is \emph{entirely rational}. \mE\ So the family/sheet
rates \emph{do} translate --- the rational angles $2\cos(\pi/2){=}0$ ($|\Z_2|$) and $2\cos(\pi/3){=}1$ ($\Nfam$)
are the same atoms that build the recovery rate $(2/3)^6=(|\Z_2|/\Nfam)^{2\Nfam}\in\mathbb Q$ --- but the carrier
$5$-fold (golden, $\mathbb Q(\sqrt5)$) has \emph{no} rational dynamic image (a field obstruction: $2/3$ is neither
an adjacency nor a lazy-update eigenvalue, and the heat-kernel time $\Delta/(2-\phi)\approx6.37$ is not clean).
\mE\ The carrier is therefore \emph{static-only} --- $\gcar{=}5$ builds the lattice and spectrum, while
$\{|\Z_2|,\Nfam\}{=}\{2,3\}$ drive the relaxation. The only object holding both is the order-$30$ Coxeter element,
$30=\gcar\,(2\Nfam)=5\cdot6$ (the golden in the order-$5$ factor, the recovery exponent $6{=}2\Nfam$ in the
order-$6$ factor); correspondingly the $\phi$-quasicrystal $E_8=H_4{+}\phi H_4$ carries the carrier but not the
(rational, $3$-fold) cusp-weight family dynamics \mC.

\par\smallskip\noindent\textbf{Couple or factorize? A cyclotomic answer}
(\veri{v315\_coxeter\_coupling.py}). The order-$30$ Coxeter element does \emph{not} dynamically entangle the two
sectors --- as a cyclic group it \emph{factorizes}, $\Z/30=\Z/5\times\Z/6$ (a direct product, decoupled, in line
with the field obstruction above). It couples them instead at the \emph{field} level: $\mathbb Q(\zeta_{30})$ is the
compositum of the carrier field $\mathbb Q(\zeta_5)$ (which contains the golden $\sqrt5=\phi+1/\phi$) and the family
field $\mathbb Q(\zeta_3)$, with degree $[\mathbb Q(\zeta_{30}){:}\mathbb Q]=8=$ rank $E_8$. \mE\ Its Galois group is
exactly $(\Z/5)^\times\times(\Z/3)^\times=\Z/4\times\Z/2=\mu_4\times\Z_2$ (order $8$, maximal element order $4$,
\emph{not} cyclic) --- so the $\mu_4$ seam clock \emph{is} the Galois group of the carrier's golden field, and the
$\Z_2$ sheet that of the family field. The carrier ($\sqrt5$, $5$-fold) and the family ($\zeta_3$, $3$-fold) are
thus conjugate sub-data of \emph{one} cyclotomic field: the static$\leftrightarrow$dynamic bridge is
\emph{arithmetic} (Galois), not a dynamical generator \mC.

\par\smallskip\noindent\textbf{Does the Galois action reach the physics? The CP phases}
(\veri{v316\_galois\_readout.py}). It does --- through the \emph{family} factor. The hexagonal CP unit
$\rho=\zeta_6=\zeta_{30}^5$ is exactly the order-$6=2\Nfam$ factor $c^5$ above (the carrier factor being
$\zeta_5=\zeta_{30}^6$), so the leading CP phases are its cyclotomic arguments:
$\delta_{\mathrm{CKM}}^{\mathrm{lead}}=\pi/3=\arg\zeta_6$ and $\delta_{\mathrm{PMNS}}=4\pi/3=\arg\zeta_6^{\,4}$ (the
power $4=|\mu_4|$), with $\zeta_6^4=-\zeta_6$ the $\Z_2$ sheet flip and $\mu_6=\mu_3\times\mu_2$ (family $\times$
sheet) \mC. \mE\ Crucially the family Galois factor itself is physical: $\mathrm{Gal}(\mathbb Q(\zeta_6)/\mathbb Q)
=\Z_2$, $\zeta_6\mapsto\zeta_6^5=\bar\zeta_6$, \emph{is} CP conjugation $\delta\mapsto-\delta$. The carrier factor
$\zeta_5$ (the golden $\sqrt5$) carries no phase --- $\mu_4$ stays static --- and the \emph{magnitudes}
($\sin^2\theta_{12}=\tfrac13-\phiz/2$, with $\phiz$ transcendental) are not cyclotomic at all. So the
$\mu_4\times\Z_2$ coupling touches the readouts exactly in the phase sector the field split predicts: the CP phases
are family-factor cyclotomic data, the magnitudes the analytic seed \mE.

\par\smallskip\noindent\textbf{Are the three generations themselves a Galois orbit?}
(\veri{v317\_galois\_family.py}). The same arithmetic governs the family count, not only the phase. The three cusp
weights $\{0,\tfrac13,\tfrac23\}$ (with $\mathrm{Spec}(Q_+)=3\alpha+1=\{1,2,3\}$, the $A_3$ exponents) are, under
$w\mapsto e^{2\pi i w}$, the cube roots of unity $\{1,\zeta_3,\zeta_3^2\}$, and the cyclic family symmetry
$w\mapsto w+\tfrac13$ is a \emph{transitive} $\mu_3$ orbit --- the three generations \emph{are} the cube-root orbit.
\mE\ The family Galois $\Z_2$ ($\zeta_3\mapsto\zeta_3^2=\bar\zeta_3$, the same CP conjugation as above) refines this
to $1{+}2$: it \emph{fixes} the rational weight $w=0$ and \emph{swaps} $w=\tfrac13\leftrightarrow\tfrac23$. And the
fixed (rational) generation is precisely the recovery \emph{attractor} --- $w=0$ carries the transfer eigenvalue $1$
(the stationary ``law'', \S\ref{sec:attractor}), while the Galois-conjugate pair carries the decaying modes
$(2/3)^6,(1/3)^6$ (the hierarchy). \mE\ So the family structure is the same arithmetic as the CP phase: a $\mu_3$
orbit, Galois-refined $1{+}2$, the fixed point the selected law; only the mass \emph{magnitudes} remain the analytic
seed \mC.
